\chapter{Marco teórico del Método de los Elementos Finitos en elasticidad plana}
\label{cap:marco_teorico}

El Método de los Elementos Finitos (MEF) es el procedimiento numérico de referencia para aproximar la solución de problemas de mecánica de medios continuos formulados como ecuaciones diferenciales en derivadas parciales. Su solidez teórica está documentada en los textos clásicos de la disciplina \autocite{zienkiewicz2013fem,bathe2014fem,cook2002concepts}.

Este capítulo tiene dos partes. La primera sitúa el trabajo: los antecedentes y el estado del arte del software para la enseñanza del MEF, que diagnostican la brecha y fijan los requisitos de la herramienta, y los principios que orientan la exposición del procedimiento de cálculo. La segunda establece los fundamentos matemáticos que sustentan el motor, restringidos a la elasticidad lineal plana con elementos isoparamétricos cuadriláteros de cuatro y de nueve nodos. Esos fundamentos siguen el orden del propio procedimiento de cálculo: de la formulación variacional a la recuperación de tensiones, pasando por la calidad de malla y por los criterios de verificación y validación.

Las obras de referencia del método —Zienkiewicz y Taylor, Bathe, Hughes, Cook y colaboradores, Strang y Fix— se citan en las ediciones consultadas, anteriores en más de diez años a este trabajo, por su carácter de fuentes clásicas. La formulación isoparamétrica, la cuadratura de Gauss, el tratamiento del bloqueo y la teoría de aproximación que aquí se exponen pertenecen al núcleo consolidado del método, y estas ediciones son las que el autor tuvo en mano para verificar cada página citada. La literatura reciente se reserva para la enseñanza del método, para la verificación y validación y para las bibliotecas numéricas empleadas.

\section{Antecedentes y estado del arte}
\label{sec:antecedentes}

El MEF se enseña tradicionalmente combinando el desarrollo analítico en pizarra con el uso de paquetes comerciales de propósito general. Esa práctica deja al estudiante entre dos extremos: las herramientas profesionales están optimizadas para producir resultados, no para mostrar el procedimiento intermedio, mientras que el desarrollo manual no escala más allá de uno o dos elementos. Sobre esa tensión, que el planteamiento del problema expone y la \autoref{sec:fundamentos-pedagogicos} retoma, la investigación educativa propone el software específicamente didáctico como puente \autocite{bishay2020teaching,lee2015interactive,perezsantiago2023fem}. Un estudio de consenso entre expertos de la industria y de la academia señala, además, que estos cursos tienden a privilegiar las destrezas de operación sobre los conceptos fundamentales \autocite[p.~1159]{perezsantiago2023fem}.

\paragraph{Alcance de esta revisión.} Los antecedentes que siguen proceden de una revisión no exhaustiva. Se partió de la literatura sobre enseñanza del MEF publicada en revistas de educación en ingeniería, se rastrearon las referencias de los trabajos hallados y se examinaron los programas que esas publicaciones documentan, con términos de búsqueda en español y en inglés relativos al software educativo para elementos finitos y al análisis de sólidos bidimensionales. El criterio de inclusión fue que el recurso se aplicara a la enseñanza del método y que su funcionamiento estuviera documentado en una publicación accesible. Las dos categorías de propósito general —las bibliotecas de programación y el software comercial— se incluyen como términos de contraste, no como recursos didácticos. De ahí que las afirmaciones de esta sección se acoten a los recursos revisados y no a la totalidad de los existentes.

Los antecedentes de software con vocación didáctica pueden agruparse en tres familias. La primera reúne las herramientas centradas en el análisis matricial de estructuras de barras, pórticos y cerchas planas. Bishay propone un constructor y un analizador de cerchas con los que el estudiante arma estructuras de barras y observa el ensamblaje y la resolución del sistema, complementados con pequeños proyectos que los propios estudiantes formulan \autocite[pp.~5-6]{bishay2020teaching}, en la línea de programas previos como ED-Beams y ED-Frames \autocite[pp.~243-244]{suarez1998edelas2d}. Estas herramientas no abordan el continuo elástico bidimensional.

La segunda familia la integran los simuladores que profundizan en un eslabón aislado del método. Lee desarrolla simuladores interactivos del procesamiento de las ecuaciones de elementos finitos \autocite{lee2015interactive} y de la visualización de los modos propios de la matriz de rigidez: modos de sólido rígido, modos espurios de energía nula y bloqueo por cortante, que no son modos de vibración \autocite{lee2015eigenmodes}. Ambos son funciones educativas incorporadas a VisualFEA, un programa comercial de análisis por elementos finitos.

La tercera reúne las herramientas de cálculo de propósito general, que no fueron concebidas como material didáctico y que dejan el procedimiento fuera del alcance visual del estudiante, aunque por motivos opuestos. Por un lado, las bibliotecas de programación, como CALFEM o FEniCS: en ellas el cálculo es enteramente inspeccionable, pero en el código fuente y no en una representación gráfica, de modo que exigen una competencia de programación previa. Por otro, los paquetes comerciales de análisis estructural, como SAP2000, ANSYS o Abaqus, que operan como cajas negras: entregan resultados sin exponer las matrices intermedias que los produjeron. Para estas dos categorías, la comparación de la \autoref{tab:comparativa} se restringe a los atributos que se desprenden de su propia naturaleza —presencia o ausencia de interfaz gráfica, régimen de licencia, requisitos de ejecución y exposición del cálculo intermedio— y no a juicios sobre la calidad didáctica de cada producto; los atributos que exigirían examinar cada producto en particular se dejan sin evaluar y se marcan con una raya. VisualFEA es el producto comercial con orientación didáctica declarada que con más frecuencia aparece en las publicaciones consultadas, y es el programa en el que están incorporados los simuladores de Lee. Al ser cerrado y de pago no se lo pudo examinar de primera mano, de modo que lo consignado procede de las publicaciones de sus propios autores; se lo incluye para no presentar el vacío como más amplio de lo que esta revisión permite sostener.

El antecedente más directo de este trabajo es ED-Elas2D, un programa educativo desarrollado en la Universitat Politècnica de Catalunya para el adiestramiento asistido por computadora en el análisis de sólidos bidimensionales por el MEF \autocite{suarez1998edelas2d}. ED-Elas2D comparte con EduFEM el dominio —elasticidad plana en tensión y en deformación plana—, la organización en pre-proceso, proceso y post-proceso, y una biblioteca de elementos que incluye los cuadriláteros Q4, Q8 y Q9. También, al menos en parte, el idioma: se desarrolló en una institución española y las capturas de pantalla de su publicación muestran una interfaz en castellano, aunque el trabajo no declara el idioma del programa. Su existencia confirma la pertinencia de un entorno educativo de esta clase y ayuda a delimitar el aporte del presente trabajo, que no reside en la lengua sino en la concepción y en la tecnología.

Frente a aquel programa de finales de la década de 1990 —software académico cuya licencia y disponibilidad la publicación no documenta, ligado a la plataforma y a las posibilidades de su época—, EduFEM se distribuye como software libre y de código abierto, expone de forma interactiva cada matriz del cálculo ($\bm{N}$, $\bm{J}$, $\bm{B}$, $\bm{D}$ y $\bm{k}_e$) sobre el elemento que el estudiante selecciona en el lienzo, genera de manera automática una memoria de cálculo con sustitución numérica paso a paso e incorpora una batería de verificación y validación documentada. La diferencia más pertinente para este trabajo está en la respuesta. Los propios autores de ED-Elas2D reconocen que su post-proceso no explica conceptos clave como el suavizado nodal de tensiones y deformaciones, y que es una herramienta meramente descriptiva que muestra los resultados principales del análisis \autocite[p.~253]{suarez1998edelas2d}. Es decir: el antecedente expone el procedimiento, pero devuelve la respuesta como lo hace el software comercial. EduFEM cubre ese hueco con la sonda —la consulta puntual que devuelve el estado de tensión en el punto que el usuario señala sobre el modelo—, que ofrece los valores en modo crudo y suavizado, con las reacciones contrastadas con la carga, el círculo de Mohr y el diagnóstico de la memoria de cálculo.

La \autoref{tab:comparativa} sitúa a EduFEM frente a estas categorías de recursos según los atributos relevantes para la enseñanza del continuo elástico bidimensional: los que describen el recurso —idioma, dimensión, elementos, licencia y requisitos de ejecución— y los que distinguen a este trabajo —la exposición del cálculo paso a paso, la cobertura del procedimiento completo, la memoria de cálculo automática, la posibilidad de observar los fenómenos numéricos y la verificación y validación publicadas junto con la herramienta—.

{\footnotesize
\setlength{\tabcolsep}{2pt}
\renewcommand{\arraystretch}{1.0}
\begin{longtable}{@{}>{\raggedright\arraybackslash}p{2.35cm}>{\raggedright\arraybackslash}p{2.35cm}*{5}{>{\raggedright\arraybackslash}p{2.06cm}}@{}}
\caption[Comparación de recursos para la enseñanza del MEF frente a EduFEM.]{Comparación de recursos para la enseñanza del MEF frente a EduFEM. Cada fila es un atributo relevante para enseñar el continuo elástico bidimensional; cada columna, un recurso o una categoría de recursos. La columna de EduFEM se lee como la propuesta de este trabajo frente al estado del arte.}\label{tab:comparativa}\\
\toprule
\textbf{Atributo} & \textbf{EduFEM (este trabajo)} & \textbf{ED-Elas2D} \autocite{suarez1998edelas2d} & \textbf{Analizador de cerchas} \autocite{bishay2020teaching} & \textbf{VisualFEA (simuladores)} \autocite{lee2015interactive,lee2015eigenmodes} & \textbf{Bibliotecas de programación} (CALFEM, FEniCS) & \textbf{Software comercial} (SAP2000 \autocite{csi2017sap2000}, ANSYS, Abaqus) \\
\midrule
\endfirsthead
\caption[]{Comparación de recursos para la enseñanza del MEF frente a EduFEM (continuación).}\\
\toprule
\textbf{Atributo} & \textbf{EduFEM (este trabajo)} & \textbf{ED-Elas2D} & \textbf{Analizador de cerchas} & \textbf{VisualFEA (simuladores)} & \textbf{Bibliotecas de programación} & \textbf{Software comercial} \\
\midrule
\endhead
\bottomrule
\endfoot
Idioma & Español & No consta (las capturas de la publicación muestran una interfaz en castellano) & Inglés & Inglés & Inglés & Inglés (localización parcial) \\
Dimensión & 2D & 2D & 2D/3D (barras) & 2D/3D & 2D/3D & 2D/3D \\
Elementos & Cuadriláteros Q4 y Q9 & Triángulos y cuadriláteros (Q4/Q8/Q9) & Barras (cercha) & Continuo, placas, láminas y sólidos & Diversos & Amplios \\
Licencia & Libre, código abierto & No consta & No declarada (código MATLAB) & Comercial & Libre & Comercial \\
Requisitos de ejecución & Instalador autocontenido (Windows) & Windows (década de 1990) & Requiere MATLAB & Requiere VisualFEA & Entorno de programación & Instalación y licencia del fabricante \\
Cálculo paso a paso & Alta: ocho módulos (M0--M7) sobre el modelo del usuario & Media (pre/proc/post) & Media-alta (código construido paso a paso) & Media-alta, en un eslabón & Visible en el código, sin visualización guiada & Baja (caja negra) \\
Procedimiento completo & Sí & Sí & Parcial (solo barras) & Parcial (un eslabón) & Parcial (por programación) & No (no didáctico) \\
Memoria de cálculo automática & Sí: PDF con sustitución numérica, trazable al modelo del usuario & No consta & No consta & Valores intermedios en pantalla, no como documento & No (la programa el usuario) & Informes de resultados, sin el desarrollo matricial \\
Fenómenos numéricos observables & Distorsión y $\det\bm{J}$ (M0, M2); sub-integración y modos espurios (M5); bloqueo por cortante Q4/Q9 (Cook); condicionamiento (memoria) & No consta & No consta & Modos espurios y bloqueo por cortante, vía modos propios de $\bm{K}$ & Por programación, sin visualización guiada & Sin exposición didáctica \\
Lectura de la respuesta & Sonda con tensiones crudas y suavizadas, reacciones frente a la carga, círculo de Mohr, diagnóstico en la memoria & Contornos y valores numéricos; sin explicación del suavizado (p.~253) & Reacciones y tensiones de barras & Valores intermedios en pantalla & Por programación & Contornos y tablas, sin exposición didáctica \\
Verificación y validación publicada con la herramienta & Sí: soluciones manufacturadas, Timoshenko y Cook (\autoref{cap:resultados}) & No consta & No consta & No consta & --- & --- \\
\end{longtable}}
{\singlespacing\notaapa{«Procedimiento completo»: el recurso recorre el procedimiento de cálculo del MEF de extremo a extremo, del mapeo isoparamétrico a la recuperación de tensiones. «No consta»: las publicaciones consultadas no documentan el atributo, lo que no implica que el recurso carezca de él. La raya marca atributos no evaluados para las categorías de propósito general, según el criterio enunciado en el texto. La comparación no es simétrica en su base de evidencia: la columna de EduFEM procede de la verificación directa del programa desarrollado en este trabajo, y las demás, de las publicaciones citadas en cada encabezado. En la fila «Cálculo paso a paso», la graduación se apoya en el rasgo que cada celda indica entre paréntesis: alta, el recurso muestra las matrices intermedias sobre el modelo que construye el usuario; media-alta, las muestra en un solo eslabón del método o a través del código; media, expone la secuencia de pre-proceso, proceso y post-proceso sin las matrices; baja, no las expone. Glosas de los códigos y símbolos de las celdas: M0 a M7 son los ocho módulos educativos de EduFEM (\autoref{sec:educacion}); $\det\bm{J}$ es el determinante del Jacobiano, indicador de la distorsión del elemento (\autoref{sec:jacobiano}); los modos espurios son los modos de energía nula que aparecen al integrar con menos puntos de Gauss de los necesarios (\autoref{sec:locking}); el condicionamiento mide la sensibilidad de la solución a los errores de redondeo; las soluciones manufacturadas, la viga de Timoshenko y la membrana de Cook son los tres casos de estudio del \autoref{cap:resultados}.}}

El contraste deja a la vista un vacío. Entre los recursos revisados, ninguno reúne de forma simultánea los ocho atributos siguientes:

\begin{enumerate}
    \item interfaz íntegramente en español;
    \item licencia libre y código abierto;
    \item dominio del continuo elástico plano;
    \item exposición interactiva del cálculo paso a paso sobre el modelo que construye el usuario;
    \item cobertura del procedimiento completo del MEF, del mapeo isoparamétrico a la recuperación de tensiones;
    \item memoria de cálculo trazable a ese modelo;
    \item observación guiada de la distorsión de malla, la sub-integración, el bloqueo por cortante y el condicionamiento;
    \item verificación y validación publicadas junto con la herramienta.
\end{enumerate}

Solo los simuladores de VisualFEA comparten uno de estos atributos: hacen visibles los modos espurios y el bloqueo a través de los modos propios de la matriz de rigidez, aunque en un producto comercial y en inglés. Salvo el antecedente directo de ED-Elas2D, la mayoría de los recursos revisados se concentra en un eslabón aislado del método o se limita a estructuras de barras \autocite[p.~17]{bishay2020teaching}, y ni siquiera ED-Elas2D, condicionado por la tecnología de su época, reúne ese conjunto.

Este diagnóstico cumple el primer objetivo específico y fija los cinco requisitos con que EduFEM se propone cubrir la brecha, que la \autoref{sec:requisitos} retoma uno por uno:

\begin{enumerate}
    \item el procedimiento de cálculo completo a la vista sobre el modelo del usuario;
    \item una memoria de cálculo reproducible a mano;
    \item una exactitud verificada y validada, publicada junto con la herramienta;
    \item la cobertura del contenido que los expertos exigen;
    \item un software libre y en español.
\end{enumerate}

\section{Principios que orientan la exposición del procedimiento}
\label{sec:fundamentos-pedagogicos}

Exponer el procedimiento de cálculo admite muchas formas, y no todas son igualmente útiles. Esta sección fija los cuatro principios con que se tomaron las decisiones de diseño del \autoref{cap:diseno}. La agrupación en cuatro principios es propia de este trabajo: ninguna de las fuentes citadas la enuncia como tal. Cada principio se apoya en la literatura sobre enseñanza del MEF que se indica en su párrafo, y describe una propiedad que el software debe tener, no un efecto sobre quien lo usa.

El primero es la \emph{visibilidad del procedimiento}. Los textos clásicos presentan la formulación matricial con rigor, pero el desarrollo manual de un problema no escala más allá de uno o dos elementos, y los paquetes de análisis entregan resultados sin revelar el procedimiento; el software didáctico se propone como puente entre ambos extremos, integrando la función educativa dentro del propio flujo de análisis \autocites[p.~157]{lee2015interactive}[p.~872]{lee2015eigenmodes}[p.~2]{bishay2020teaching}. De aquí se deriva que EduFEM exponga cada matriz sobre el modelo, en las mismas fases en que el análisis profesional las produce, y que muestre la respuesta junto con el procedimiento que la genera.

El segundo es la \emph{interactividad con vínculo al modelo}. Lee vincula en ambos sentidos cada entrada de la matriz de rigidez con la geometría del elemento y actualiza la visualización ante cada acción del usuario; reporta, con una evaluación cualitativa sobre sus propios cursos y sin grupo de comparación, una comprensión más intuitiva del procesamiento de las ecuaciones \autocites[pp.~158, 163, 166, 168]{lee2015interactive}[p.~874]{lee2015eigenmodes}. De aquí proceden las capas superpuestas de EduFEM, el conmutador entre la fórmula y su valor numérico, y la conversión entre elementos con refinamiento de la malla: operaciones que exigen manipular el modelo y no solo mirarlo.

El tercero es la \emph{construcción paso a paso}. Bishay hace que los estudiantes construyan por etapas las herramientas de cálculo —rigidez elemental, ensamblaje, solución, reacciones y tensiones— y mide una mejora entre cohortes, que atribuye principalmente a que los propios estudiantes formulan sus problemas y solo en parte a las herramientas computacionales \autocite[pp.~4-5, 13, 17]{bishay2020teaching}. De aquí proceden el orden de los módulos M1 a M7, que sigue la secuencia con que la literatura de referencia presenta la formulación isoparamétrica \autocite{bathe2014fem,cook2002concepts}, y la memoria de cálculo, que reconstruye el desarrollo con sustitución numérica sobre el problema del propio usuario, rotula cada etapa y lo entrega en un solo documento.

El cuarto es la \emph{conexión entre el fundamento y la operación}. El consenso de expertos ya citado encuentra que la formación tiende a privilegiar las destrezas de operación sobre los conceptos fundamentales, y reclama que el egresado sepa planificar, verificar y validar un análisis, no solo ejecutarlo \autocite[pp.~1159, 1168]{perezsantiago2023fem}. Las destrezas de mallado, post-proceso, conceptos fundamentales y verificación y validación que ese consenso enumera \autocite[p.~1163]{perezsantiago2023fem} son las que el \autoref{cap:diseno} adopta como universo de contenido. De aquí que la herramienta incorpore la verificación y la validación como parte de lo que expone, y no solo como garantía de su corrección.

Estos cuatro principios orientan el diseño y explican por qué EduFEM expone lo que expone y del modo en que lo hace. No son evidencia de que la herramienta produzca un efecto de aprendizaje: ese efecto no se mide en este trabajo, y las fuentes citadas lo reportan de otras herramientas y en condiciones que sus propios autores acotan. Lo que sí permiten es comprobar que las decisiones de diseño no son arbitrarias, y el \autoref{cap:diseno} registra, para cada una, el principio del que procede.

\section{Fundamentos del MEF y formulación variacional en elasticidad lineal}
\label{sec:formulacion-debil}

El análisis matricial de pórticos, materia corriente en la formación de ingeniería civil, ofrece el punto de partida más cercano. En él, cada barra aporta una matriz de rigidez que liga las fuerzas de sus extremos con los desplazamientos de sus nudos; esas matrices elementales se ensamblan en una matriz global y el sistema $\bm{K}\,\bm{u} = \bm{F}$ devuelve los desplazamientos de todos los nudos de la estructura. El MEF hace lo mismo con un sólido continuo: lo divide en elementos, obtiene la rigidez de cada uno, los ensambla en la misma matriz global y resuelve el mismo sistema. La diferencia está en el origen de esa rigidez elemental. La barra tiene solución exacta conocida, y la teoría de vigas entrega su matriz en forma cerrada; un trozo de chapa no la tiene, de modo que su rigidez no se escribe directamente, sino que se obtiene integrando sobre el área del elemento. Todo lo que sigue en este capítulo sirve a ese cálculo: cómo se interpola el desplazamiento dentro del elemento, cómo se derivan de él las deformaciones, cómo se integra numéricamente sobre el área, cómo se ensambla el sistema y cómo se recuperan después las tensiones.

Sea un dominio bidimensional $\Omega \subset \mathbb{R}^2$ con frontera $\Gamma = \Gamma_u \cup \Gamma_t$, donde $\Gamma_u$ es la porción con desplazamientos prescritos (condiciones de Dirichlet, es decir, los apoyos) y $\Gamma_t$ la porción con tracciones impuestas (condiciones de Neumann, es decir, las cargas de borde). El problema de equilibrio en elasticidad lineal estática se enuncia, en su forma fuerte, como

\begin{equation}
\nabla \cdot \bm{\sigma} + \bm{b} = \bm{0} \quad \text{en } \Omega,
\label{eq:equilibrio-fuerte}
\end{equation}

donde $\bm{\sigma}$ es el tensor de tensiones y $\bm{b}$ la fuerza por unidad de volumen. La forma fuerte exige continuidad de segundo orden del campo de desplazamientos $\bm{u}$, requisito difícil de satisfacer de manera global. El MEF parte por ello de la \emph{forma débil} o variacional, un enunciado equivalente del equilibrio que pide menos suavidad al campo de desplazamientos. Se obtiene multiplicando la \autoref{eq:equilibrio-fuerte} por un campo de desplazamientos virtuales $\delta\bm{u}$ admisible (nulo sobre $\Gamma_u$, para que respete los apoyos) e integrando por partes mediante el teorema de la divergencia. Esto conduce al principio de los trabajos virtuales \autocites[pp.~69-71]{zienkiewicz2013fem}[pp.~77-79]{hughes2000fem}[pp.~612-613]{reddy2006introduction}:

\begin{equation}
\int_\Omega \delta\bm{\varepsilon}^{T} \bm{\sigma} \, d\Omega
= \int_\Omega \delta\bm{u}^{T} \bm{b} \, d\Omega
+ \int_{\Gamma_t} \delta\bm{u}^{T} \bm{t} \, d\Gamma,
\label{eq:trabajos-virtuales}
\end{equation}

con $\delta\bm{\varepsilon}$ las deformaciones virtuales y $\bm{t}$ las tracciones de borde. La formulación débil rebaja el requisito de continuidad a primer orden, lo que habilita aproximaciones polinómicas a trozos. El campo continuo se discretiza interpolando los valores nodales mediante funciones de forma: dentro de cada elemento, $\bm{u} \approx \bm{N}\,\bm{u}_e$, donde $\bm{N}$ agrupa las funciones de forma y $\bm{u}_e$ los desplazamientos de los nodos del elemento. Sustituyendo esta aproximación en la \autoref{eq:trabajos-virtuales} y aprovechando la arbitrariedad de $\delta\bm{u}$ se obtiene el sistema algebraico de equilibrio

\begin{equation}
\bm{K}\,\bm{u} = \bm{F},
\label{eq:sistema-global}
\end{equation}

donde $\bm{K}$ es la matriz de rigidez global, $\bm{u}$ el vector de desplazamientos nodales y $\bm{F}$ el vector de cargas. El resto del capítulo desarrolla cada pieza de la \autoref{eq:sistema-global}.

\section{Elasticidad plana: tensión plana y deformación plana}
\label{sec:elasticidad-plana}

La elasticidad bidimensional admite dos idealizaciones de un sólido tridimensional \autocites[pp.~11-12]{timoshenko1970elasticity}[pp.~79, 94]{cook2002concepts}. La \emph{tensión plana} modela cuerpos delgados cargados en su plano, en los que la tensión normal al plano es despreciable ($\sigma_z = 0$); es el caso de una chapa o membrana de poco espesor. La \emph{deformación plana} modela cuerpos prismáticos largos en los que la deformación longitudinal se anula ($\varepsilon_z = 0$); corresponde, por ejemplo, a una presa o un tubo largo donde una rebanada representa todo el sólido.

En ambos casos las componentes relevantes se agrupan en notación de Voigt como $\bm{\sigma} = [\sigma_x,\, \sigma_y,\, \tau_{xy}]^T$ y $\bm{\varepsilon} = [\varepsilon_x,\, \varepsilon_y,\, \gamma_{xy}]^T$, ligadas por la ley constitutiva $\bm{\sigma} = \bm{D}\,\bm{\varepsilon}$. Para tensión plana la matriz constitutiva es

\begin{equation}
\bm{D}_{\text{TP}} = \frac{E}{1-\nu^{2}}
\begin{bmatrix}
1 & \nu & 0 \\
\nu & 1 & 0 \\
0 & 0 & \dfrac{1-\nu}{2}
\end{bmatrix},
\label{eq:D-tension-plana}
\end{equation}

y para deformación plana

\begin{equation}
\bm{D}_{\text{DP}} = \frac{E}{(1+\nu)(1-2\nu)}
\begin{bmatrix}
1-\nu & \nu & 0 \\
\nu & 1-\nu & 0 \\
0 & 0 & \dfrac{1-2\nu}{2}
\end{bmatrix},
\label{eq:D-deformacion-plana}
\end{equation}

donde $E$ es el módulo de elasticidad y $\nu$ el coeficiente de Poisson. El cambio en el denominador refleja el acoplamiento adicional que introduce la restricción $\varepsilon_z = 0$. La deformación plana es sensible a la incompresibilidad: cuando $\nu \to 0{,}5$ el factor $(1-2\nu)$ tiende a cero y $\bm{D}_{\text{DP}}$ se degrada, fenómeno asociado al bloqueo volumétrico (\emph{volumetric locking}) que se analiza en la \autoref{sec:locking}.

\section{Elementos isoparamétricos Q4 y Q9}
\label{sec:isoparametricos}

EduFEM emplea elementos isoparamétricos cuadriláteros: isoparamétrico significa que las mismas funciones de forma describen la geometría del elemento y el campo de desplazamientos dentro de él, con iguales parámetros para una cosa y para la otra \autocites[p.~339]{bathe2014fem}[p.~187]{onate2009structural}. Las funciones se definen sobre un elemento de referencia, el cuadrado natural $(\xi, \eta) \in [-1, 1]^2$, y un mapeo isoparamétrico las lleva a la geometría física del elemento real de la malla.

El elemento \textbf{Q4} es un cuadrilátero de cuatro nodos con interpolación bilineal. Numerados sus vértices en sentido antihorario, las funciones de forma son

\begin{equation}
\begin{aligned}
N_1 &= \tfrac{1}{4}(1-\xi)(1-\eta), & N_2 &= \tfrac{1}{4}(1+\xi)(1-\eta), \\
N_3 &= \tfrac{1}{4}(1+\xi)(1+\eta), & N_4 &= \tfrac{1}{4}(1-\xi)(1+\eta).
\end{aligned}
\label{eq:N-q4}
\end{equation}

El elemento \textbf{Q9} añade cuatro nodos medios de arista y un nodo central, alcanzando interpolación bicuadrática como producto tensorial de polinomios de Lagrange unidimensionales. Sus funciones de forma para vértices, lados y centro son, respectivamente,

\begin{equation}
\begin{aligned}
N_1 &= \tfrac{1}{4}\,\xi(\xi-1)\,\eta(\eta-1), & N_2 &= \tfrac{1}{4}\,\xi(\xi+1)\,\eta(\eta-1), \\
N_3 &= \tfrac{1}{4}\,\xi(\xi+1)\,\eta(\eta+1), & N_4 &= \tfrac{1}{4}\,\xi(\xi-1)\,\eta(\eta+1), \\
N_5 &= \tfrac{1}{2}(1-\xi^2)\,\eta(\eta-1), & N_6 &= \tfrac{1}{2}\,\xi(\xi+1)(1-\eta^2), \\
N_7 &= \tfrac{1}{2}(1-\xi^2)\,\eta(\eta+1), & N_8 &= \tfrac{1}{2}\,\xi(\xi-1)(1-\eta^2), \\
N_9 &= (1-\xi^2)(1-\eta^2). &&
\end{aligned}
\label{eq:N-q9}
\end{equation}

Todas satisfacen la propiedad de Kronecker $N_i(\xi_j, \eta_j) = \delta_{ij}$, es decir, cada función vale uno en su propio nodo y cero en los demás, de modo que el valor interpolado en un nodo es exactamente el de ese nodo. Y todas forman una partición de la unidad: en cualquier punto del elemento las funciones suman uno, lo que garantiza por construcción que el elemento reproduzca sin error un desplazamiento de cuerpo rígido. El mapeo isoparamétrico interpola la posición física de cualquier punto interior a partir de las coordenadas nodales $(x_i, y_i)$:

\begin{equation}
x(\xi, \eta) = \sum_{i=1}^{n} N_i(\xi, \eta)\,x_i, \qquad
y(\xi, \eta) = \sum_{i=1}^{n} N_i(\xi, \eta)\,y_i,
\label{eq:mapeo-iso}
\end{equation}

con $n = 4$ para Q4 y $n = 9$ para Q9. El elemento Q9, de orden superior, captura gradientes de tensión más complejos y tolera mejor la distorsión geométrica de la malla \autocites[pp.~192, 208]{hughes2000fem}[pp.~476-477]{bathe2014fem}, a costa de un mayor número de grados de libertad (GDL): 8 GDL en Q4 frente a 18 GDL en Q9. La \autoref{fig:isoparametricos} ilustra ambos elementos en el cuadrado natural y su mapeo al espacio físico.

\begin{figure}[htbp]
\centering
\includegraphics[width=\linewidth]{fig_isoparametricos}
\caption[Elementos isoparamétricos Q4 y Q9 y su mapeo al espacio físico.]{Elementos isoparamétricos Q4 y Q9. A la izquierda, el elemento de referencia: un cuadrado de lado dos en coordenadas naturales $(\xi,\eta)$ entre $-1$ y $1$, con los nodos numerados en sentido antihorario desde la esquina inferior izquierda. A la derecha, el elemento real de la malla, al que el mapeo isoparamétrico lleva cada punto del cuadrado. Esta numeración es la que emplean el motor de cálculo y la matriz de extrapolación de la \autoref{eq:extrap-q4}.}
\label{fig:isoparametricos}
\end{figure}

\section{Jacobiano del mapeo}
\label{sec:jacobiano}

El mapeo de la \autoref{eq:mapeo-iso} relaciona los dos espacios mediante la matriz Jacobiana, el cuadro de derivadas que registra cuánto se estira y se gira el elemento al pasar del cuadrado de referencia a su forma real y que convierte derivadas respecto de las coordenadas naturales en derivadas respecto de las físicas:

\begin{equation}
\bm{J} =
\begin{bmatrix}
\dfrac{\partial x}{\partial \xi} & \dfrac{\partial y}{\partial \xi} \\[6pt]
\dfrac{\partial x}{\partial \eta} & \dfrac{\partial y}{\partial \eta}
\end{bmatrix}
= \frac{\partial \bm{N}}{\partial(\xi,\eta)}\,\bm{X}_e,
\label{eq:jacobiano}
\end{equation}

donde $\partial\bm{N}/\partial(\xi,\eta)$ es la matriz $2\times n$ cuyas filas son las derivadas $\partial N_i/\partial\xi$ y $\partial N_i/\partial\eta$, y $\bm{X}_e$ ($n\times 2$) reúne las coordenadas nodales del elemento. Esta es la convención de Bathe \autocite[p.~345]{bathe2014fem}, que es también la que implementa el motor. Con ella, cada fila de $\bm{J}$ es el vector tangente a una línea de coordenada natural ($\partial\bm{x}/\partial\xi$ y $\partial\bm{x}/\partial\eta$), y la regla de la cadena de la \autoref{eq:fisica-natural} emplea $\bm{J}^{-1}$ directamente. La convención transpuesta, también frecuente en la literatura, exigiría en su lugar $\bm{J}^{-\mathsf{T}}$. El determinante del Jacobiano mide el factor local de cambio de área entre ambos espacios, $dx\,dy = |\det\bm{J}|\,d\xi\,d\eta$: es el número por el que hay que multiplicar un área del cuadrado de referencia para obtener la del elemento real, y por eso resulta central en la integración numérica. Para una matriz $2\times2$ se calcula de forma explícita,

\begin{equation}
\det\bm{J} = J_{11}J_{22} - J_{12}J_{21},
\label{eq:detJ}
\end{equation}

y su inversa se obtiene por cofactores, evitando rutinas genéricas de inversión:

\begin{equation}
\bm{J}^{-1} = \frac{1}{\det\bm{J}}
\begin{bmatrix}
J_{22} & -J_{12} \\
-J_{21} & J_{11}
\end{bmatrix}.
\label{eq:invJ}
\end{equation}

Exigir $\det\bm{J} > 0$ en todo el elemento, con el determinante de la \autoref{eq:detJ}, garantiza por construcción que el mapeo sea biyectivo y el elemento válido \autocites[pp.~346-347]{bathe2014fem}[p.~236]{cook2002concepts}. Un determinante nulo o negativo indica un elemento degenerado o invertido, y un programa de cálculo debe detectarlo antes de resolver; el umbral concreto con que lo hace EduFEM se fija en la \autoref{sec:motor}.

\section{Matriz de deformación-desplazamiento B}
\label{sec:matriz-b}

Las deformaciones se obtienen derivando el campo de desplazamientos interpolado. La relación cinemática en elasticidad plana es

\begin{equation}
\bm{\varepsilon} =
\begin{bmatrix} \varepsilon_x \\ \varepsilon_y \\ \gamma_{xy} \end{bmatrix}
=
\begin{bmatrix}
\dfrac{\partial u}{\partial x} \\[4pt]
\dfrac{\partial v}{\partial y} \\[4pt]
\dfrac{\partial u}{\partial y} + \dfrac{\partial v}{\partial x}
\end{bmatrix},
\label{eq:cinematica}
\end{equation}

que en forma matricial se escribe $\bm{\varepsilon} = \bm{B}\,\bm{u}_e$, donde $\bm{u}_e$ reúne los desplazamientos de los nodos del elemento. La matriz $\bm{B}$, de dimensión $3 \times 2n$, es la que convierte esos desplazamientos nodales en deformaciones, y contiene para ello las derivadas físicas de las funciones de forma:

\begin{equation}
\bm{B} =
\begin{bmatrix}
\dfrac{\partial N_1}{\partial x} & 0 & \dfrac{\partial N_2}{\partial x} & 0 & \cdots \\[6pt]
0 & \dfrac{\partial N_1}{\partial y} & 0 & \dfrac{\partial N_2}{\partial y} & \cdots \\[6pt]
\dfrac{\partial N_1}{\partial y} & \dfrac{\partial N_1}{\partial x} & \dfrac{\partial N_2}{\partial y} & \dfrac{\partial N_2}{\partial x} & \cdots
\end{bmatrix}.
\label{eq:matriz-B}
\end{equation}

Las funciones de forma se definen en coordenadas naturales, pero $\bm{B}$ exige sus derivadas físicas $\partial N_i/\partial x$ y $\partial N_i/\partial y$. El puente lo provee la inversa del Jacobiano de la \autoref{eq:invJ}, aplicada por columna a cada función de forma según la regla de la cadena:

\begin{equation}
\begin{bmatrix}
\dfrac{\partial N_i}{\partial x} \\[6pt]
\dfrac{\partial N_i}{\partial y}
\end{bmatrix}
= \bm{J}^{-1}
\begin{bmatrix}
\dfrac{\partial N_i}{\partial \xi} \\[6pt]
\dfrac{\partial N_i}{\partial \eta}
\end{bmatrix}.
\label{eq:fisica-natural}
\end{equation}

Esta relación vectorial es la conexión esencial entre el espacio natural, donde las derivadas se conocen analíticamente, y el espacio físico, donde se formula la deformación.

\section{Matriz de rigidez elemental e integración numérica de Gauss}
\label{sec:rigidez-gauss}

La matriz de rigidez de un elemento integra el producto $\bm{B}^T\bm{D}\bm{B}$ sobre su dominio físico, con $t$ el espesor:

\begin{equation}
\bm{k}_e = \int_{\Omega_e} \bm{B}^T \bm{D}\, \bm{B}\, t \; dA
= \int_{-1}^{1}\!\!\int_{-1}^{1} \bm{B}^T(\xi,\eta)\,\bm{D}\,\bm{B}(\xi,\eta)\, t\, |\det\bm{J}|\; d\xi\, d\eta.
\label{eq:ke-integral}
\end{equation}

El integrando no es polinómico: como $\bm{B}$ depende de $\bm{J}^{-1}$, que es racional en $(\xi,\eta)$, la integración algebraica exacta deja de ser practicable salvo en los elementos más simples \autocites[p.~147]{zienkiewicz2013fem}[p.~191]{hughes2000fem}. Por ello se aproxima mediante cuadratura de Gauss-Legendre, que evalúa el integrando en puntos óptimos y lo pondera:

\begin{equation}
\bm{k}_e \approx \sum_{g=1}^{n_g} w_g\; \bm{B}^T(\xi_g,\eta_g)\,\bm{D}\,\bm{B}(\xi_g,\eta_g)\, t\, |\det\bm{J}(\xi_g,\eta_g)|,
\label{eq:ke-gauss}
\end{equation}

donde $g$ recorre los $n_g$ puntos de cuadratura del elemento, y los puntos $(\xi_g,\eta_g)$ y los pesos $w_g$ surgen del producto tensorial de la regla unidimensional. La regla de $m$ puntos integra exactamente polinomios de grado $2m-1$ \autocites[pp.~141-142]{hughes2000fem}[p.~161]{zienkiewicz2013fem}. La integración completa de esta familia es la cuadratura $2\times2$ para Q4 (cuatro puntos en $\xi = \pm 1/\sqrt{3}$) y $3\times3$ para Q9 (puntos en $\xi = \pm\sqrt{3/5},\,0$ con pesos $5/9,\,8/9,\,5/9$), órdenes que integran exactamente la rigidez de los respectivos elementos no distorsionados \autocites[pp.~221-222]{hughes2000fem}[p.~470]{bathe2014fem}; es la que adopta el motor (\autoref{sec:motor}).

El orden de integración no es arbitrario. Una \emph{subintegración} puede dejar la matriz de rigidez con rango deficiente y admitir modos de deformación de energía nula, los modos espurios (\emph{hourglass}) que se analizan en la \autoref{sec:locking}. Una \emph{sobreintegración}, por ejemplo $3\times3$ en Q4, no aporta precisión apreciable y solo encarece el cálculo: el integrando solo es polinómico en los elementos no distorsionados, y en los distorsionados la diferencia es marginal. La elección $2\times2$/$3\times3$ equilibra exactitud y estabilidad para esta familia de elementos.

\section{Fenómenos numéricos en elementos de bajo orden: bloqueo y modos espurios}
\label{sec:locking}

La discretización con elementos de bajo orden introduce fenómenos numéricos que no provienen de la física del problema, sino de las limitaciones de la interpolación, y cuya comprensión resulta esencial para interpretar correctamente los resultados \autocite{cook2002concepts,hughes2000fem,bathe2014fem}. Tres son particularmente relevantes para los cuadriláteros que emplea EduFEM y reaparecen, ya cuantificados, en el \autoref{cap:resultados}.

El \textbf{bloqueo por cortante} (\emph{shear-locking}) afecta al elemento bilineal Q4 en estados dominados por flexión. Bajo flexión pura, una sección debería curvarse sin desarrollar deformación de corte; sin embargo, el campo de desplazamientos bilineal del Q4 no puede representar esa curvatura sin generar al mismo tiempo deformaciones de corte espurias, ausentes en la solución exacta. Esas deformaciones parásitas absorben energía y vuelven al elemento artificialmente rígido, de modo que subestima los desplazamientos y converge con lentitud al refinar la malla \autocites[pp.~98-99]{cook2002concepts}[p.~243]{hughes2000fem}. El elemento Q9, cuya interpolación bicuadrática sí reproduce el modo de curvatura, está prácticamente libre de este bloqueo \autocites[pp.~235-236]{cook2002concepts}[p.~243]{hughes2000fem}. El \autoref{cap:resultados} lo cuantifica.

El \textbf{bloqueo volumétrico} (\emph{volumetric locking}) aparece en deformación plana cuando el material se aproxima a la incompresibilidad ($\nu \to 0{,}5$). Como anticipó la \autoref{sec:elasticidad-plana}, el factor $(1-2\nu)$ del denominador de $\bm{D}_{\text{DP}}$ tiende a cero, y la restricción de incompresibilidad sobre-condiciona el elemento: la malla no dispone de grados de libertad suficientes para satisfacer la condición $\varepsilon_x + \varepsilon_y \to 0$ en cada punto sin rigidizarse en exceso. El resultado es, de nuevo, una respuesta artificialmente rígida, acompañada a menudo de oscilaciones espurias en el campo de tensiones \autocites[pp.~192, 208]{hughes2000fem}[pp.~282-283]{bathe2014fem}.

Los \textbf{modos espurios} (\emph{hourglass}) son el fenómeno inverso: en lugar de rigidizar el elemento, lo ablandan. Como se anticipó en la \autoref{sec:rigidez-gauss}, surgen al \emph{subintegrar} la rigidez, por ejemplo evaluando el Q4 con un único punto de Gauss. La matriz elemental queda entonces con rango deficiente y admite modos de deformación de energía nula: formas de deformación que el elemento adopta sin oponer resistencia alguna, porque en el único punto donde se mide la deformación esta resulta nula, y que dibujan sobre la malla el patrón de reloj de arena que da nombre al fenómeno \autocites[p.~223]{cook2002concepts}[pp.~472-473]{bathe2014fem}. Al no costar energía, esos patrones no físicos pueden crecer sin control y contaminar la solución. EduFEM los evita empleando la integración completa de cada elemento ($2\times2$ para Q4 y $3\times3$ para Q9).

Existen técnicas correctivas para los dos primeros bloqueos. Una es la integración reducida selectiva (SRI), que integra con menos puntos de Gauss únicamente la parte de la rigidez responsable del bloqueo y deja el resto con la cuadratura completa \autocites[p.~221]{hughes2000fem}[pp.~476-477]{bathe2014fem}. Otra es la formulación $\bar{\bm{B}}$ (\emph{B-bar}), que promedia sobre el elemento esa misma parte del operador de deformación en lugar de evaluarla punto por punto \autocite[pp.~232-234]{hughes2000fem}. EduFEM no las implementa: ofrece en su lugar el elemento Q9, que evita el bloqueo por cortante por la vía del mayor orden de interpolación, aunque no el volumétrico, frente al cual el bicuadrático con integración completa también requiere integración selectiva \autocite[p.~222]{hughes2000fem}. Además, conserva el comportamiento del Q4 sin corregirlo, de modo que el bloqueo quede visible en los resultados que el programa presenta en lugar de quedar compensado dentro del cálculo. La incorporación de esas técnicas correctivas se plantea como línea de extensión futura entre las recomendaciones del trabajo.

\section{Ensamblaje, restricciones y cargas}
\label{sec:ensamblaje}

Las matrices elementales de la \autoref{eq:ke-gauss} y los vectores de carga de cada elemento se ensamblan en el sistema global mediante el mapeo de cada GDL local a su índice global. La matriz de rigidez global se obtiene con el operador de ensamblaje, que aquí se denota $\Ensamblaje$ y que no es otra cosa que la suma de cada aporte elemental en las filas y columnas que le corresponden según ese mapeo:

\begin{equation}
\bm{K} = \Ensamblaje_{e} \bm{k}_e, \qquad
\bm{F} = \sum_e \bm{f}_e^{\text{vol}} + \sum_j \bm{F}_j + \sum_\Gamma \bm{F}_\Gamma,
\label{eq:ensamblaje}
\end{equation}

donde $e$ recorre los elementos de la malla, $j$ los nodos cargados y $\Gamma$ las aristas con carga de superficie, y donde un nodo compartido por varios elementos acumula los aportes de todos ellos. Para mallas grandes $\bm{K}$ es dispersa (mayoritariamente nula), por lo que EduFEM la construye en formato de coordenadas (COO) y la convierte a filas comprimidas (CSR) para el solucionador, apoyándose en el módulo de matrices dispersas de SciPy \autocite[p.~265]{virtanen2020scipy}.

El vector de cargas reúne tres aportes. Las \textbf{cargas nodales} $\bm{F}_j$ se aplican directamente en los GDL del nodo $j$. Las \textbf{cargas volumétricas} $\bm{f}_e^{\text{vol}}$, como la gravedad, se distribuyen consistentemente integrando $\int_{\Omega_e} \bm{N}^T \bm{b}\, t\, dA$ por cuadratura de Gauss. Las \textbf{cargas superficiales} $\bm{F}_\Gamma$ sobre una arista se convierten en fuerzas nodales equivalentes ponderando la carga distribuida con las funciones de forma. Para una arista lineal de Q4 de longitud $L$ con carga $q$ que varía linealmente entre los valores $q_1$ y $q_2$ de sus extremos, el reparto exacto es $F_1 = \tfrac{L}{6}(2q_1 + q_2)$ y $F_2 = \tfrac{L}{6}(q_1 + 2q_2)$, que con $q$ constante se reduce a $qL/2$ en cada nodo. Para la arista cuadrática de Q9 con $q$ constante la distribución resulta $\tfrac{qL}{6},\,\tfrac{4qL}{6},\,\tfrac{qL}{6}$, concentrando la mayor parte en el nodo medio \autocite[p.~113]{cook2002concepts}.

Las restricciones de Dirichlet se imponen partiendo el sistema en GDL libres ($f$) y restringidos ($r$):

\begin{equation}
\begin{bmatrix} \bm{K}_{ff} & \bm{K}_{fr} \\ \bm{K}_{rf} & \bm{K}_{rr} \end{bmatrix}
\begin{bmatrix} \bm{u}_f \\ \bm{u}_r \end{bmatrix}
=
\begin{bmatrix} \bm{F}_f \\ \bm{F}_r \end{bmatrix}.
\label{eq:particion}
\end{equation}

Para \emph{restricciones homogéneas} ($\bm{u}_r = \bm{0}$) basta resolver el bloque libre de la \autoref{eq:particion}, $\bm{K}_{ff}\,\bm{u}_f = \bm{F}_f$. Para \emph{restricciones no homogéneas} ($\bm{u}_r \neq \bm{0}$), necesarias en la verificación con soluciones manufacturadas, el sistema reducido incorpora el término de sustitución estática

\begin{equation}
\bm{K}_{ff}\,\bm{u}_f = \bm{F}_f - \bm{K}_{fr}\,\bm{u}_r,
\label{eq:dirichlet-no-homogeneo}
\end{equation}

que preserva la simetría del bloque libre. Un modelo bien restringido hace que $\bm{K}_{ff}$ sea definida positiva, es decir, que toda deformación distinta de cero almacene energía, condición que asegura una solución única y, por tanto, que la matriz no sea singular \autocite[pp.~726-727]{bathe2014fem}. Bien restringido significa aquí al menos tres GDL, dispuestos de modo que ninguna combinación de ellos permita las dos traslaciones ni la rotación del sólido como cuerpo rígido: tres apoyos paralelos o tres apoyos concurrentes no bastan, igual que no bastan en el análisis de un pórtico.

\section{Resolución del sistema y recuperación de tensiones}
\label{sec:resolucion-tensiones}

El sistema reducido de la \autoref{eq:dirichlet-no-homogeneo} se resuelve por factorización LU directa sobre la matriz dispersa, a través de la interfaz que SciPy expone sobre la biblioteca SuperLU \autocite[p.~264]{virtanen2020scipy}. El orden en que se eliminan las incógnitas condiciona el llenado (\emph{fill-in}) de los factores, de modo que la factorización dispersa se precede de un reordenamiento de las ecuaciones; la elección concreta del motor se documenta en la \autoref{sec:motor}. Como $\bm{K}_{ff}$ es simétrica y definida positiva, admite además una factorización de Cholesky, $\bm{K}_{ff} = \bm{L}\bm{L}^{\top}$, que almacena un solo triángulo y no requiere pivoteo; el motor emplea la factorización LU general de SuperLU, y la de Cholesky se reserva como línea de trabajo (\emph{Recomendaciones}). Conocido el vector $\bm{u}$, las reacciones en los GDL restringidos se recuperan como $\bm{R} = \bm{K}_{rf}\,\bm{u}_f + \bm{K}_{rr}\,\bm{u}_r - \bm{F}_r$.

Las tensiones se calculan primero en los puntos de Gauss, evaluando $\bm{\sigma}_g = \bm{D}\,\bm{B}_g\,\bm{u}_e$ en cada punto $g$. A partir de las componentes cartesianas se obtienen las tensiones principales \autocite[pp.~12-14]{timoshenko1970elasticity} y la equivalente de von Mises, la medida escalar con que habitualmente se resume un estado tensional biaxial \autocite[p.~117]{cook2002concepts}:

\begin{equation}
\begin{aligned}
\sigma_{1,2} &= \frac{\sigma_x + \sigma_y}{2} \pm \sqrt{\left(\frac{\sigma_x - \sigma_y}{2}\right)^2 + \tau_{xy}^2}, \\
\sigma_{\mathrm{VM}} &= \sqrt{\tfrac{1}{2}\left[(\sigma_1-\sigma_2)^2 + (\sigma_2-\sigma_z)^2 + (\sigma_z-\sigma_1)^2\right]}, \\
\sigma_z &=
\begin{cases}
0 & \text{tensión plana,}\\
\nu\,(\sigma_x+\sigma_y) & \text{deformación plana.}
\end{cases}
\end{aligned}
\label{eq:principales-vm}
\end{equation}

En tensión plana $\sigma_z = 0$ y la \autoref{eq:principales-vm} se reduce a la forma plana $\sigma_{\mathrm{VM}} = \sqrt{\sigma_1^2 - \sigma_1\sigma_2 + \sigma_2^2}$. En deformación plana la tensión fuera del plano no es nula y es la tercera tensión principal. Omitirla, simplificación a la que invita el carácter bidimensional del planteo, subestima la tensión equivalente; por eso EduFEM la incorpora en todas las rutas de cálculo del programa —tensiones en los puntos de Gauss, valores extrapolados a los nodos, promedio nodal, consulta puntual sobre el modelo y memoria de cálculo (\autoref{sec:motor})—, con el $\nu$ del material de cada elemento.

Las tensiones no son igualmente exactas en todo el elemento: existen puntos de muestreo en los que convergen con un orden superior al del resto, propiedad que se conoce como superconvergencia y que en los elementos de esta familia corresponde a los puntos de Gauss de una regla de orden inferior al de la integración completa. En el Q9 son los cuatro puntos de la regla $2\times2$, y no los nueve de la regla $3\times3$ con que se integra su rigidez \autocite[pp.~230-231]{cook2002concepts}. EduFEM evalúa las tensiones en los puntos de la cuadratura de rigidez de cada elemento y las lleva a los nodos mediante una extrapolación, es decir, prolongando hacia las esquinas los valores medidos en el interior, con una matriz precalculada: $\bm{\sigma}_{\text{nodos}} = \bm{E}\,\bm{\sigma}_{\text{Gauss}}$, en lugar de evaluar $\bm{B}$ directamente en los nodos. Conviene advertirlo una vez: $\bm{E}$, en negrita, designa la matriz de extrapolación y no el módulo de elasticidad $E$ de la \autoref{eq:D-tension-plana}, que se escribe sin negrita. Para el Q4 esa elección muestrea en los cuatro puntos de la regla $2\times2$, aunque Cook sitúa el punto superconvergente del elemento bilineal en su centro \autocite[p.~230]{cook2002concepts}: es una decisión de implementación que conserva la variación del campo dentro del elemento. Para el Q9 privilegia que la matriz $\bm{M}$ definida a continuación sea cuadrada con los nueve nodos del elemento, a costa de no muestrear en los puntos óptimos.

El costo esperable de esa decisión es doble: las tensiones nodales no heredan la precisión superior de los puntos óptimos, y la extrapolación amplifica las diferencias entre puntos de Gauss, pues los pesos de la matriz del Q4 llegan a $1{,}866$ y $-0{,}5$, como se ve más abajo en la \autoref{eq:extrap-q4}. A cambio, una única ruta de cálculo sirve a los dos tipos de elemento y conserva la variación del campo dentro de cada uno. El orden de convergencia que resulta de esta secuencia de muestreo, extrapolación y promediado no se toma de la literatura: se mide en el \autoref{cap:resultados} con el método de las soluciones manufacturadas.

La matriz de extrapolación $\bm{E}$ se deduce de la propia interpolación isoparamétrica. Dentro del elemento, el campo de tensiones se aproxima con las mismas funciones de forma que la geometría, las de la \autoref{eq:N-q4} y la \autoref{eq:N-q9}, de modo que el valor en un punto de Gauss es la interpolación de los valores nodales,

\begin{equation}
\bm{\sigma}(\xi_g,\eta_g) = \sum_{i=1}^{n} N_i(\xi_g,\eta_g)\,\bm{\sigma}_{\text{nodo},i}.
\label{eq:interp-tension}
\end{equation}

Evaluando la \autoref{eq:interp-tension} en los puntos de Gauss del elemento se obtiene el sistema $\bm{\sigma}_{\text{Gauss}} = \bm{M}\,\bm{\sigma}_{\text{nodos}}$, donde $M_{gi} = N_i(\xi_g,\eta_g)$ reúne las funciones de forma del nodo $i$ evaluadas en cada punto de Gauss $g$. Cuando el número de puntos de Gauss iguala al de nodos del elemento —cuadratura $2\times2$ para Q4 y $3\times3$ para Q9—, la matriz $\bm{M}$ es cuadrada e invertible, y la extrapolación es su inversa,

\begin{equation}
\bm{\sigma}_{\text{nodos}} = \bm{M}^{-1}\,\bm{\sigma}_{\text{Gauss}} = \bm{E}\,\bm{\sigma}_{\text{Gauss}}.
\label{eq:extrap-inversa}
\end{equation}

Para el elemento Q4 esta inversa admite forma cerrada. Como los cuatro puntos de Gauss se ubican en $(\pm 1/\sqrt{3},\,\pm 1/\sqrt{3})$ y los nodos en $(\pm 1,\,\pm 1)$, extrapolar equivale a evaluar la interpolación bilineal en coordenadas escaladas por $\sqrt{3}$, con entradas de la forma $\tfrac{1}{4}(1\pm\sqrt{3})(1\pm\sqrt{3})$. Con los puntos de Gauss numerados en el orden que produce el motor —$(-,-)$, $(-,+)$, $(+,-)$, $(+,+)$ en $(\xi,\eta)$— y los nodos en sentido antihorario desde $(-1,-1)$, resulta

\begin{equation}
\bm{E}_{\text{Q4}} =
\begin{bmatrix}
a & b & b & c \\
b & c & a & b \\
c & b & b & a \\
b & a & c & b
\end{bmatrix},
\qquad
a = 1+\frac{\sqrt{3}}{2},\quad
b = -\frac{1}{2},\quad
c = 1-\frac{\sqrt{3}}{2},
\label{eq:extrap-q4}
\end{equation}

con $a \approx 1{,}866$ y $c \approx 0{,}134$: cada nodo recibe el peso $a$ del punto de Gauss más próximo, $c$ del diagonalmente opuesto y $b$ de los otros dos. La numeración no es un detalle menor: permutar los puntos de Gauss permuta las columnas de $\bm{E}$, y una matriz escrita para otra numeración corrompe las tensiones nodales sin alterar los desplazamientos, de modo que una verificación basada solo en desplazamientos no lo detecta. Por eso el motor construye $\bm{E}$ invirtiendo numéricamente $\bm{M}$ con sus propios puntos de Gauss, para Q4 igual que para Q9, y una prueba de regresión comprueba que $\bm{E}$ reproduce los campos bilineales (o bicuadráticos) de prueba con un error del orden de la precisión de máquina, que es la verificación numérica de la identidad algebraica de la que $\bm{E}$ se deduce. Para el Q9, $\bm{M}$ es de $9\times9$ y su inversa no admite forma cerrada simple. La extrapolación se aplica a las componentes cartesianas $(\sigma_x,\sigma_y,\tau_{xy})$; las tensiones principales y la de von Mises, al ser funciones no lineales de ellas, se recalculan a partir de las componentes ya extrapoladas en lugar de extrapolarse de forma directa: es el criterio que adopta EduFEM para que cada valor nodal sea consistente con sus propias componentes cartesianas.

Finalmente, como un nodo compartido recibe valores distintos desde cada elemento adyacente, se aplica un \emph{promediado nodal}

\begin{equation}
\bm{\sigma}_{\text{nodo}} = \frac{1}{n_e}\sum_{e=1}^{n_e} \bm{\sigma}_{\text{nodo}}^{(e)},
\label{eq:promediado}
\end{equation}

que suaviza el campo y produce un contorno continuo. El salto entre los valores sin promediar es, además, un indicador útil del error de discretización —es el principio en que se apoyan los estimadores de error por recuperación \autocite[pp.~476-477]{zienkiewicz2013fem}—: cuanto mayor la dispersión entre elementos, más gruesa es la malla en esa zona.

\section{Calidad de malla}
\label{sec:calidad-malla}

La precisión del resultado depende de la calidad geométrica de los elementos. EduFEM evalúa dos métricas complementarias, tomadas de la biblioteca de calidad geométrica Verdict \autocite[p.~3]{stimpson2007verdict}, que miden propiedades independientes entre sí: una los ángulos del elemento y otra sus proporciones. Un rectángulo muy alargado, por ejemplo, tiene ángulos rectos impecables y proporciones inadmisibles. El \textbf{Jacobiano escalado} mide la distorsión angular normalizando el determinante del Jacobiano por las normas de sus filas, que son los vectores tangentes a las líneas $\xi$ y $\eta$,

\begin{equation}
q_{SJ} = \min_g \frac{\det\bm{J}_g}{\|\bm{J}_{1,:}\|_g\,\|\bm{J}_{2,:}\|_g},
\label{eq:scaled-jacobian}
\end{equation}

cantidad que equivale al seno del ángulo local entre ambas tangentes y vive en el rango $[-1, 1]$: vale $1$ cuando las dos tangentes son perpendiculares, indica baja calidad cerca de $0$ y delata un elemento invertido cuando $q_{SJ} \le 0$ \autocite[p.~51]{stimpson2007verdict}. El comprobador de salud (\autoref{sec:pre-proceso}) y la memoria de cálculo la evalúan en los puntos de Gauss del elemento. El módulo introductorio de calidad de malla (\autoref{sec:educacion}) emplea la forma por vértices de Verdict, el seno del ángulo interior en cada esquina tomando el peor de los cuatro, que es la forma de la métrica que el módulo presenta paso a paso, calculable a mano con el ángulo de cada esquina. La \textbf{compacidad} o \emph{stretch} mide la proporción del elemento relacionando su arista más corta $L_{\min}$ con su diagonal más larga $D_{\max}$,

\begin{equation}
q_C = \sqrt{2}\,\frac{L_{\min}}{D_{\max}} \in [0, 1],
\label{eq:stretch}
\end{equation}

que vale $1$ para un cuadrado perfecto y disminuye al alargarse o aplastarse el elemento \autocite[p.~57]{stimpson2007verdict}. El subíndice $C$ recuerda que se trata de la compacidad, y evita confundir el símbolo con los nombres Q4 y Q9 de los elementos. De la documentación de Verdict provienen los mínimos admisibles del cuadrilátero, $q_{SJ} \ge 0{,}30$ y $q_C \ge 0{,}25$ \autocite[pp.~51, 57]{stimpson2007verdict}; los cortes con que el programa los endurece y clasifica los elementos son decisiones de implementación, documentadas en la \autoref{sec:pre-proceso}. El programa asigna además a cada elemento un estado global de calidad —buena, aceptable o mala— combinando cuatro métricas. Las dos primeras salen del propio Jacobiano: el $q_{SJ}$ de la \autoref{eq:scaled-jacobian}, evaluado en los puntos de Gauss, y la razón $R_J = \min_g \det\bm{J}_g / \max_g \det\bm{J}_g$ entre el menor y el mayor de sus determinantes. Las dos restantes se construyen sobre una descomposición bilineal del cuadrilátero, equivalente a la que Verdict emplea para sus métricas de forma \autocite[pp.~35-36]{stimpson2007verdict}, que reduce las coordenadas de las cuatro esquinas $\bm{r}_1,\dots,\bm{r}_4$ a tres vectores característicos:

\begin{equation}
\begin{aligned}
\bm{X}_1 &= \tfrac{1}{4}\left(-\bm{r}_1 + \bm{r}_2 + \bm{r}_3 - \bm{r}_4\right), \qquad
\bm{X}_2 = \tfrac{1}{4}\left(-\bm{r}_1 - \bm{r}_2 + \bm{r}_3 + \bm{r}_4\right), \\[2pt]
\bm{X}_3 &= \tfrac{1}{4}\left(\bm{r}_1 - \bm{r}_2 + \bm{r}_3 - \bm{r}_4\right), \\[4pt]
AR &= \frac{\max(\|\bm{X}_1\|, \|\bm{X}_2\|)}{\min(\|\bm{X}_1\|, \|\bm{X}_2\|)}, \qquad
T_R = \max\left(\frac{|\bm{X}_3 \cdot \bm{X}_1|}{\|\bm{X}_1\|^2},\; \frac{|\bm{X}_3 \cdot \bm{X}_2|}{\|\bm{X}_2\|^2}\right).
\end{aligned}
\label{eq:descomposicion-bilineal}
\end{equation}

Los dos primeros vectores unen los puntos medios de las aristas opuestas en cada dirección natural y miden la extensión del elemento; el tercero recoge su parte propiamente bilineal. De ellos se derivan la \textbf{relación de aspecto} $AR$, que vale $1$ en un cuadrado o un rombo y crece al alargarse el elemento, y el \textbf{estrechamiento} (\emph{taper}) $T_R$, que es nulo en todo paralelogramo y crece con la trapecialidad: precisamente el modo de distorsión ante el que el Q4 es más sensible, como muestra el caso de Cook del \autoref{cap:resultados}. En el Q9 la cuarta métrica no es el estrechamiento, sino la desviación de los nodos de arista y central respecto de su posición ideal sobre el elemento macro. Los cortes con que se asigna el estado global no son umbrales bibliográficos sino convención del programa, y se enumeran en la \autoref{sec:pre-proceso}.

\section{Verificación y validación}
\label{sec:verificacion-validacion}

Un programa de cálculo debe distinguir entre \emph{verificación}, que comprueba que las ecuaciones se resuelven correctamente, y \emph{validación}, que comprueba que las ecuaciones representan la realidad física \autocite[pp.~13-14]{oberkampf2010vv}. Este trabajo aplica ambas estrategias en la acepción que se precisa a continuación: verificación en sentido estricto, y validación entendida como contraste con referencias externas de exactitud conocida, no como campaña experimental. En su acepción estricta, en efecto, la validación exige el contraste con mediciones experimentales, y este trabajo no realiza campaña experimental propia. Sus referencias externas son una solución analítica de la teoría de la elasticidad y un modelo equivalente resuelto en un programa comercial, de modo que el contraste corresponde a verificación de solución y a comparación código a código; esta última, según la propia literatura, no sustituye a la verificación rigurosa del código \autocite[pp.~173-174]{oberkampf2010vv}.

En lo que sigue, y en el resto del documento, «validación» designa ese contraste con referencias externas de exactitud conocida, es decir, con la solución analítica y con el modelo comercial. Es la acepción corriente en la práctica del MEF y es distinta de la validación experimental en sentido estricto, que no se realiza. El documento emplea además «cobertura de contenido» para el contraste documental del \autoref{cap:resultados}, que no es validación en el sentido de esta sección, y «comprobador de salud» para la revisión automática de los datos de entrada, que es el nombre de un componente del programa. La sigla V\&V, de uso habitual en la disciplina, abrevia «verificación y validación» en ese doble sentido.

La verificación de código se realiza con el \textbf{Método de las Soluciones Manufacturadas} (MMS): se elige a priori un campo de desplazamientos exacto $\bm{u}_M$, se sustituye en las ecuaciones de equilibrio para deducir la fuerza de cuerpo $\bm{b} = -\nabla\cdot\bm{\sigma}(\bm{u}_M)$ que lo produce, y se inyectan esa carga y las restricciones de Dirichlet no homogéneas correspondientes \autocites[pp.~21, 26]{salari2000mms}[pp.~219-222]{oberkampf2010vv}.

Para comparar la solución del programa con la exacta no basta mirar un punto: hace falta una medida única del error en toda la pieza. La norma $L^2$ del desplazamiento es esa medida, una media cuadrática de la diferencia entre el desplazamiento calculado y el exacto, ponderada por el área, de modo que resume en un solo número el desajuste promedio en todo el dominio. La seminorma $H^1$ hace lo mismo con los gradientes del desplazamiento, es decir, con las deformaciones, que son las que determinan las tensiones que el ingeniero acaba leyendo; por eso converge más despacio y es la medida más exigente de las dos. Al reducir a la mitad el tamaño $h$ de los elementos, ambas medidas deben caer en una proporción fija que la teoría predice: esa proporción es la \emph{tasa de convergencia}, y comprobarla es la prueba más severa de que el programa resuelve lo que dice resolver. La proporción depende del \emph{orden} $p$ del elemento, que es el grado de los polinomios con que interpola: $p = 1$ para el Q4 bilineal y $p = 2$ para el Q9 bicuadrático. El orden fija cuánto cuesta ganar exactitud: con tasa $2$, esa reducción a la mitad divide el error por cuatro; con tasa $3$, por ocho.

Con esa lectura, el error de la solución numérica frente a la exacta se cuantifica con la norma $L^2$ del desplazamiento y la seminorma $H^1$:

\begin{equation}
\begin{aligned}
\|\bm{u}_h - \bm{u}\|_{L^2} &= \left( \int_\Omega \|\bm{u}_h - \bm{u}\|^2 \, d\Omega \right)^{1/2}, \\
|\bm{u}_h - \bm{u}|_{H^1} &= \left( \int_\Omega \|\nabla\bm{u}_h - \nabla\bm{u}\|^2 \, d\Omega \right)^{1/2}.
\end{aligned}
\label{eq:normas-error}
\end{equation}

En la \autoref{eq:normas-error}, $\bm{u}$ es la solución exacta manufacturada y $\bm{u}_h$ la calculada sobre una malla de tamaño característico $h$; las barras dobles corresponden a la norma $L^2$ y las simples a la seminorma $H^1$. Más abajo aparece además $\bm{\sigma}^*_h$, donde el asterisco distingue el campo de tensiones recuperado, extrapolado a los nodos y promediado, del evaluado directamente en los puntos de Gauss.

Al refinar la malla, estas medidas deben decrecer con las tasas asintóticas teóricas ($\mathcal{O}(h^{p+1})$ en $L^2$ y $\mathcal{O}(h^p)$ en la seminorma $H^1$, con $p$ el orden polinómico del elemento) que establece la teoría de aproximación del MEF \autocites[pp.~106-107]{strang2008analysis}[pp.~189-190]{hughes2000fem}. Una tasa observada inferior a la teórica delata un defecto de implementación; alcanzarla es la evidencia más exigente de que las rutinas ejercitadas por esa medida están correctamente implementadas, sin que ello descarte errores en las rutas que la medida no recorre.

Las integrales de error se evalúan con un punto de Gauss más por dirección que la rigidez, es decir, $3\times3$ para Q4 y $4\times4$ para Q9, para no subestimar artificialmente el error. Es una precaución de implementación de los propios guiones de verificación: el integrando de la norma es el cuadrado de una diferencia entre dos campos, de grado más alto que el integrando de la rigidez, y una cuadratura ajustada a esta última lo muestrea justo en los puntos donde el error de discretización tiende a anularse, de modo que el error de la cuadratura se confunde con el que se quiere medir y degrada el orden observado.

Como esas medidas dependen solo del campo de desplazamientos, no ejercitan la secuencia de recuperación de tensiones, esto es, la evaluación en los puntos de Gauss, la extrapolación a los nodos y el promediado. Por ello se mide además la norma $L^2$ del campo de tensiones recuperado $\bm{\sigma}^*_h$ frente a $\bm{\sigma}(\bm{u}_M)=\bm{D}\,\bm{\varepsilon}(\bm{u}_M)$, porque el error de las magnitudes derivadas debe verificarse por separado y su orden puede ser menor que el de la solución primaria \autocite[p.~29]{salari2000mms}. Además, se comprueba de forma independiente que la matriz de extrapolación reproduce los campos polinómicos de su espacio de interpolación con un error del orden de la precisión de máquina.

La validación, en la acepción fijada arriba, contrasta el motor contra casos de estudio con solución conocida. La \textbf{viga de Timoshenko} simplemente apoyada admite solución analítica de la teoría de la elasticidad \autocite[art.~21, pp.~39-43]{timoshenko1970elasticity}; el \autoref{cap:resultados} contrasta con ella la tensión normal y la flecha, y compara además contra un modelo de cáscara de SAP2000 \autocite[p.~179]{csi2017sap2000}. La \textbf{membrana de Cook} es el caso clásico con que la literatura evidencia la sensibilidad de los elementos a la distorsión geométrica y al bloqueo por cortante: en ella el elemento Q4 subestima sensiblemente el desplazamiento del extremo en las mallas gruesas, por el bloqueo por cortante descrito en la \autoref{sec:locking} \autocite[pp.~98-99]{cook2002concepts}, mientras que el Q9, de orden superior, lo evita y converge con rapidez \autocite[p.~243]{hughes2000fem}. El valor de uso extendido para ese desplazamiento, $23{,}96$, se corrobora en este trabajo con una extrapolación de Richardson propia, es decir, con el procedimiento que estima el valor exacto desconocido a partir de tres mallas sucesivas; la \autoref{sec:cook} presenta el caso completo, con su geometría, su valor de referencia y esa corroboración. Estos contrastes muestran por qué la elección del tipo de elemento condiciona la calidad del resultado.

Con lo expuesto en este capítulo queda establecida la base común del motor de cálculo, de los módulos educativos y de la memoria de cálculo: una misma formulación, las mismas matrices y los mismos criterios de verificación sirven a los tres, de modo que lo que el motor calcula, lo que los módulos muestran sobre el modelo y lo que la memoria escribe paso a paso son el mismo procedimiento leído con distinto detalle. Con ello se cumple el segundo objetivo específico. El \autoref{cap:diseno} traduce esta base en decisiones de diseño y de implementación.
